\documentclass[oneside,12pt,numbers,spanish]{ezthesis}
%% # Opciones disponibles para el documento #
%%
%% Las opciones con un (*) son las opciones predeterminadas.
%%
%% Modo de compilar:
%%   draft            - borrador con marcas de fecha y sin im'agenes
%%   draftmarks       - borrador con marcas de fecha y con im'agenes
%%   final (*)        - version final de la tesis
%%
%% Tama'no de papel:
%%   letterpaper (*)  - tama'no carta (Am'erica)
 %%   a4paper          - tama'no A4    (Europa)
%%
%% Formato de impresi'on:
%%   oneside          - hojas impresas por un solo lado
%%   twoside (*)      - hijas impresas por ambos lados
%%
%% Tama'no de letra:
%%   10pt, 11pt, o 12pt (*)
%%
%% Espaciado entre renglones:
%%   singlespace      - espacio sencillo
%%   onehalfspace (*) - espacio de 1.5
%%   doublespace      - a doble espacio
%%
%% Formato de las referencias bibliogr'aficas:
%%   numbers          - numeradas, p.e. [1]
%%   authoryear (*)   - por autor y a'no, p.e. (Newton, 1997)
%%
%% Opciones adicionales:
%%   spanish         - tesis escrita en espa'nol
%%
%% Desactivar opciones especiales:
%%   nobibtoc   - no incluir la bibiolgraf'ia en el 'Indice general
%%   nofancyhdr - no incluir "fancyhdr" para producir los encabezados
%%   nocolors   - no incluir "xcolor" para producir ligas con colores
%%   nographicx - no incluir "graphicx" para insertar gr'aficos
%%   nonatbib   - no incluir "natbib" para administrar la bibliograf'ia

%% Paquetes adicionales requeridos se pueden agregar tambi'en aqu'i.
%% Por ejemplo:
\usepackage[spanish]{babel}
\usepackage{cite}
\usepackage{natbib} 
\usepackage{subfig}
\usepackage{multirow, array}
\usepackage{graphicx} 
\usepackage{float}  


%% # Datos del documento #
%%\documentclass[12pt]{•}%% Nota que los acentos se deben escribir: \'a, \'e, \'i, etc.
%% La letra n con tilde es: \~n.

\author{Stephanie Correa}
\title{Desarrollo de una aplicaci\'on Web}
\degree{Ingeniero en Inform\'atica}
\supervisor{Tutor}
\institution{Universidad Nacional Experimental Del T\'achira}
\faculty{Vicerectorado Acad\'emico}
\department{Departamento de Ingenier\'ia en Inform\'atica}

%% # M'argenes del documento #
%% 
%% Quitar el comentario en la siguiente linea para austar los m'argenes del
%% documento. Leer la documentaci'on de "geometry" para m'as informaci'on.

\geometry{left=3cm,right=3cm,top=4cm,bottom=3cm}

%% El siguiente comando agrega ligas activas en el documento para las
%% referencias cruzadas y citas bibliogr'aficas. Tiene que ser *la 'ultima*
%% instrucci'on antes de \begin{document}.
\hyperlinking
\begin{document}

%% En esta secci'on se describe la estructura del documento de la tesis.
%% Consulta los reglamentos de tu universidad para determinar el orden
%% y la cantidad de secciones que debes de incluir.

%% # Portada de la tesis #
%% Mirar el archivo "titlepage.tex" para los detalles.
%% ## Construye tu propia portada ##
%% 
%% Una portada se conforma por una secuencia de "Blocks" que incluyen
%% piezas individuales de informaci'on. Un "Block" puede incluir, por
%% ejemplo, el t'itulo del documento, una im'agen (logotipo de la universidad),
%% el nombre del autor, nombre del supervisor, u cualquier otra pieza de
%% informaci'on.
%%
%% Cada "Block" aparece centrado horizontalmente en la p'agina y,
%% verticalmente, todos los "Blocks" se distruyen de manera uniforme 
%% a lo largo de p'agina.
%%
%% Nota tambi'en que, dentro de un mismo "Block" se pueden cortar
%% lineas usando el comando \\
%%
%% El tama'no del texto dentro de un "Block" se puede modificar usando uno de
%% los comandos:
%%   \small      \LARGE
%%   \large      \huge
%%   \Large      \Huge
%%
%% Y el tipo de letra se puede modificar usando:
%%   \bfseries - negritas
%%   \itshape  - it'alicas
%%   \scshape  - small caps
%%   \slshape  - slanted
%%   \sffamily - sans serif
%%
%% Para producir plantillas generales, la informaci'on que ha sido inclu'ida
%% en el archivo principal "tesis.tex" se puede accesar aqu'i usando:
%%   \insertauthor
%%   \inserttitle
%%   \insertsupervisor
%%   \insertinstitution
%%   \insertdegree
%%   \insertfaculty
%%   \insertdepartment
%%   \insertsubmitdate

\begin{titlepage}
\vspace*{-4cm}
  \TitleBlock{\includegraphics[scale=0.2]{unet.jpg} }  
  \TitleBlock{\small\insertinstitution}
  \TitleBlock[\small]{\insertfaculty}
  \TitleBlock[\small]{Decanato de Docencia}
  \TitleBlock[\small]{\insertdepartment}
  \TitleBlock[\small]{Trabajo de aplicaci\'on profesional}
  \TitleBlock[\small]{Proyecto especial de Grado}
   \vspace*{3cm}
    
 \TitleBlock{\scshape\inserttitle}
\vspace*{3cm}
 
 \begin{flushright}
  \TitleBlock[\small]{ Autor(es): \insertauthor} 
  \TitleBlock[\small]{ C.I.: }
  \TitleBlock[\small]{ correo@gmail.com}
  \TitleBlock[\small]{ Tutor(es): \insertsupervisor}
  \TitleBlock[\small]{ correo@gmail.com}
 \end{flushright}
  \TitleBlock{\small San Crist\'obal, \insertsubmitdate}
  
\end{titlepage}

%% Nota 1:
%% Se puede agregar un escudo o logotipo en un "Block" como:
%%   \TitleBlock{\includegraphics[height=4cm]{escudo_uni}}
%% y teniendo un archivo "escudo_uni.pdf", "escudo_uni.png" o "escudo_uni.jpg"
%% en alg'un lugar donde LaTeX lo pueda encontrar.

%% Nota 2:
%% Normalmente, el espacio entre "Blocks" se extiende de modo que el
%% contenido se reparte uniformemente sobre toda la p'agina. Este
%% comportamiento se puede modificar para mantener fijo, por ejemplo, el
%% espacio entre un par de "Blocks". Escribiendo:
%%   \TitleBlock{Bloque 1}
%%   \TitleBlock[\bigskip]{Bloque2}
%% se deja un espacio "grande" y de tama~no fijo entre el bloque 1 y 2.
%% Adem'as de \bigskip est'an tambi'en \smallskip y \medskip. Si necesitas
%% aun m'as control puedes usar tambi'en, por ejemplo, \vspace*{2cm}.




%% # Prefacios #
%% Por cada prefacio (p.e. agradecimientos, resumen, etc.) crear
%% un nuevo archivo e incluirlo aqu'i.
%% Para m'as detalles y un ejemplo mirar el archivo "gracias.tex".
%%\include{gracias}
\chapter*{}
\pagestyle{empty}
\thispagestyle{empty}
\begin{flushright}
San Crist\'obal, Mes de 2016
\end{flushright}


\begin{flushleft}
Se\~nores:
\textbf{Miembros de la Comisi\'on de Trabajo de Aplicaci\'on profesional}
\textit{Departamento de Ingenier\'ia Inform\'atica}
\end{flushleft}



	Yo, Nombre del autor, titular de la c\'edula de identidad N V12.345.678, inscrita en el per\'iodo acad\'emico semestre, estudiante del noveno semestre de Ingenier\'ia Inform\'atica, por medio de la presente, someto a consideraci\'on de la Comsi\'on de Trabajo de Aplicaci\'on Profesional de este Departamento la propuesta de Proyecto Especial de Grado titulado:\textbf{Desarrollo de una aplicaci\'on Web}, la cual se desarrollar\'a bajo la tutor\'ia de la profesor(a): Tutor.\\
	
	
	Se anexan los recaudos exigidos, de acuerdo con lo dispuesto en la normativa para el Trabajo de Aplicaci\'on Profesional de la Universidad nacional Experimental del T\'achira. la Fecha estimada de culminaci\'on del proyecto:fecha.\\
	
	
	Sin otro particular a que hacer referencia y en espera de su respuesta,\\
	
	
	Atentamente,\\
	
	
	\begin{center}
	Nombre completo del autor
	
	C.I. V12.345.678
	\end{center}
	 
\chapter*{}
\pagestyle{empty}
\thispagestyle{empty}


\begin{center}
\includegraphics[scale=0.2]{unet.jpg}\\

Universidad Nacional Experimental del T\'achira\\

Vicerrectorado Acad\'emico\\

Decanato de Docencia\\

Departamento de Ingenier\'ia Inform\'atica\\

Trabajo de Aplicaci\'on Profesional\\

Proyecto Especial de Grado\\

\end{center}

\begin{center}
\textbf{Aprobaci\'on del Tutor para presentaci\'on de la Propuesta del Proyecto Especial de Grado}
\end{center}

	Yo, Nombre y apellidos del Tutor titular de la c\'edula de identidad: 12.345.678 en mi car\'acter de Tutor(a) del Proyecto Especial de Grado titulado: Desarrollo de una aplicaci\'on Web, presentado por el bachiller: Nombre completo del autor titular de la c\'edula de identidad N V12.345.678, por medio de la presente autorizo la presentaci\'on de la Propuesta de Proyecto Especial de Grado ante los jurados designados por la comisi\'on de Trabajo de Aplicaci\'on Profesional del Departamento de Ingenier\'ia en Inform\'atica, en virtud de considerar que re\'une los requisitos establecidos en el art\'iculo 16 de las Normas para el Trabajo de Aplicaci\'on Profesional de la UNET.\\
	
	
	\begin{center}
	Nombres y apellidos del tutor\\
	
	c\'edula de identidad del tutor
	\end{center}

%% # 'Indices y listas de contenido #
%% Quitar los comentarios en las lineas siguientes para obtener listas de
%% figuras y cuadros/tablas.
\pagenumbering{roman}
\setcounter{page}{4}
\cleardoublepage
\addcontentsline{toc}{chapter}{\'Indice General}
\tableofcontents
\renewcommand{\listtablename}{\'Indice de Tablas}
\cleardoublepage
\addcontentsline{toc}{chapter}{\'Indice de Tablas} % para que aparezca en el indice de contenidos
\listoftables % indice de tablas
\cleardoublepage
\addcontentsline{toc}{chapter}{\'Indice de Figuras} % para que aparezca en el indice de contenidos
\listoffigures % indice de figuras



%% # Cap'itulos #
%% Por cada cap'itulo hay que crear un nuevo archivo e incluirlo aqu'i.
%% Mirar el archivo "intro.tex" para un ejemplo y recomendaciones para
%% escribir.
\cleardoublepage
\addcontentsline{toc}{chapter}{Introducci\'on}
\markboth{Introducci\'on}{Introducci\'on} % encabezado 

\include{intro}
\chapter{El Problema}

\section{Planteamiento y Formulaci\'on}

	lorem ipsum sdfjs fdsdj fshdbfsf dfsndjdf sd sbdf fwbefiubw bfiwbfib lorem ipsum sdfjs fdsdj fshdbfsf dfsndjdf sd sbdf fwbefiubw bfiwbfib lorem ipsum sdfjs fdsdj fshdbfsf dfsndjdf sd sbdf fwbefiubw bfiwbfib lorem ipsum sdfjs fdsdj fshdbfsf dfsndjdf sd sbdf fwbefiubw bfiwbfib lorem ipsum sdfjs fdsdj fshdbfsf dfsndjdf sd sbdf fwbefiubw bfiwbfib lorem ipsum sdfjs fdsdj fshdbfsf dfsndjdf sd sbdf fwbefiubw bfiwbfib lorem ipsum sdfjs fdsdj fshdbfsf dfsndjdf sd sbdf fwbefiubw bfiwbfib lorem ipsum sdfjs fdsdj fshdbfsf dfsndjdf sd sbdf fwbefiubw bfiwbfib

	
	lorem ipsum sdfjs fdsdj fshdbfsf dfsndjdf sd sbdf fwbefiubw bfiwbfib lorem ipsum sdfjs fdsdj fshdbfsf dfsndjdf sd sbdf fwbefiubw bfiwbfib lorem ipsum sdfjs fdsdj fshdbfsf dfsndjdf sd sbdf fwbefiubw bfiwbfib lorem ipsum sdfjs fdsdj fshdbfsf dfsndjdf sd sbdf fwbefiubw bfiwbfib lorem ipsum sdfjs fdsdj fshdbfsf dfsndjdf sd sbdf fwbefiubw bfiwbfib lorem ipsum sdfjs fdsdj fshdbfsf dfsndjdf sd sbdf fwbefiubw bfiwbfib lorem ipsum sdfjs fdsdj fshdbfsf dfsndjdf sd sbdf fwbefiubw bfiwbfib lorem ipsum sdfjs fdsdj fshdbfsf dfsndjdf sd sbdf fwbefiubw bfiwbfib

	lorem ipsum sdfjs fdsdj fshdbfsf dfsndjdf sd sbdf fwbefiubw bfiwbfib lorem ipsum sdfjs fdsdj fshdbfsf dfsndjdf sd sbdf fwbefiubw bfiwbfib lorem ipsum sdfjs fdsdj fshdbfsf dfsndjdf sd sbdf fwbefiubw bfiwbfib lorem ipsum sdfjs fdsdj fshdbfsf dfsndjdf sd sbdf fwbefiubw bfiwbfib lorem ipsum sdfjs fdsdj fshdbfsf dfsndjdf sd sbdf fwbefiubw bfiwbfib lorem ipsum sdfjs fdsdj fshdbfsf dfsndjdf sd sbdf fwbefiubw bfiwbfib lorem ipsum sdfjs fdsdj fshdbfsf dfsndjdf sd sbdf fwbefiubw bfiwbfib lorem ipsum sdfjs fdsdj fshdbfsf dfsndjdf sd sbdf fwbefiubw bfiwbfib

\section{Objetivos}
	
	
\subsection{Objetivo General}

	lorem ipsum sdfjs fdsdj fshdbfsf dfsndjdf sd sbdf fwbefiubw 

\subsection{Objetivos Espec\'ificos}
 
\begin{itemize}
\item lorem ipsum sdfjs fdsdj fshdbfsf dfsndjdf sd sbdf fwbefiubw 
\item lorem ipsum sdfjs fdsdj fshdbfsf dfsndjdf sd sbdf fwbefiubw 
\item lorem ipsum sdfjs fdsdj fshdbfsf dfsndjdf sd sbdf fwbefiubw .

\end{itemize}

\section{Justificaci\'on e Importancia}

	lorem ipsum sdfjs fdsdj fshdbfsf dfsndjdf sd sbdf fwbefiubw bfiwbfib lorem ipsum sdfjs fdsdj fshdbfsf dfsndjdf sd sbdf fwbefiubw bfiwbfib lorem ipsum sdfjs fdsdj fshdbfsf dfsndjdf sd sbdf fwbefiubw bfiwbfib lorem ipsum sdfjs fdsdj fshdbfsf dfsndjdf sd sbdf fwbefiubw bfiwbfib lorem ipsum sdfjs fdsdj fshdbfsf dfsndjdf sd sbdf fwbefiubw bfiwbfib lorem ipsum sdfjs fdsdj fshdbfsf dfsndjdf sd sbdf fwbefiubw bfiwbfib lorem ipsum sdfjs fdsdj fshdbfsf dfsndjdf sd sbdf fwbefiubw bfiwbfib lorem ipsum sdfjs fdsdj fshdbfsf dfsndjdf sd sbdf fwbefiubw bfiwbfib
	
\section{Alcance y Limitaciones}

	lorem ipsum sdfjs fdsdj fshdbfsf dfsndjdf sd sbdf fwbefiubw bfiwbfib lorem ipsum sdfjs fdsdj fshdbfsf dfsndjdf sd sbdf fwbefiubw bfiwbfib lorem ipsum sdfjs fdsdj fshdbfsf dfsndjdf sd sbdf fwbefiubw bfiwbfib lorem ipsum sdfjs fdsdj fshdbfsf dfsndjdf sd sbdf fwbefiubw bfiwbfib lorem ipsum sdfjs fdsdj fshdbfsf dfsndjdf sd sbdf fwbefiubw bfiwbfib lorem ipsum sdfjs fdsdj fshdbfsf dfsndjdf sd sbdf fwbefiubw bfiwbfib lorem ipsum sdfjs fdsdj fshdbfsf dfsndjdf sd sbdf fwbefiubw bfiwbfib lorem ipsum sdfjs fdsdj fshdbfsf dfsndjdf sd sbdf fwbefiubw bfiwbfib
\chapter{Fundamentos Te\'oricos}

\section{Antecedentes}

	lorem ipsum sdfjs fdsdj fshdbfsf dfsndjdf sd sbdf fwbefiubw bfiwbfib lorem ipsum sdfjs fdsdj fshdbfsf dfsndjdf sd sbdf fwbefiubw bfiwbfib lorem ipsum sdfjs fdsdj fshdbfsf dfsndjdf sd sbdf fwbefiubw bfiwbfib lorem ipsum sdfjs fdsdj fshdbfsf dfsndjdf sd sbdf fwbefiubw bfiwbfib lorem ipsum sdfjs fdsdj fshdbfsf dfsndjdf sd sbdf fwbefiubw bfiwbfib lorem ipsum sdfjs fdsdj fshdbfsf dfsndjdf sd sbdf fwbefiubw bfiwbfib lorem ipsum sdfjs fdsdj fshdbfsf dfsndjdf sd sbdf fwbefiubw bfiwbfib lorem ipsum sdfjs fdsdj fshdbfsf dfsndjdf sd sbdf fwbefiubw bfiwbfib

\section{Bases Te\'oricas}
	lorem ipsum sdfjs fdsdj fshdbfsf dfsndjdf sd sbdf fwbefiubw bfiwbfib lorem ipsum sdfjs fdsdj fshdbfsf dfsndjdf sd sbdf fwbefiubw bfiwbfib lorem ipsum sdfjs fdsdj fshdbfsf dfsndjdf sd sbdf fwbefiubw bfiwbfib lorem ipsum sdfjs fdsdj fshdbfsf dfsndjdf sd sbdf fwbefiubw bfiwbfib lorem ipsum sdfjs fdsdj fshdbfsf dfsndjdf sd sbdf fwbefiubw bfiwbfib lorem ipsum sdfjs fdsdj fshdbfsf dfsndjdf sd sbdf fwbefiubw bfiwbfib

\section{Definici\'on de T\'erminos}

lorem ipsum sdfjs fdsdj fshdbfsf dfsndjdf sd sbdf fwbefiubw bfiwbfib


\chapter{Fundamentos Metodol\'ogicos}

	
\section{Enfoque de la investigaci\'on}
	
	lorem ipsum sdfjs fdsdj fshdbfsf dfsndjdf sd sbdf fwbefiubw bfiwbfib lorem ipsum sdfjs fdsdj fshdbfsf dfsndjdf sd sbdf fwbefiubw bfiwbfib lorem ipsum sdfjs fdsdj fshdbfsf dfsndjdf sd sbdf fwbefiubw bfiwbfib lorem ipsum sdfjs fdsdj fshdbfsf dfsndjdf sd sbdf fwbefiubw bfiwbfib lorem ipsum sdfjs fdsdj fshdbfsf dfsndjdf sd sbdf fwbefiubw bfiwbfib lorem ipsum sdfjs fdsdj fshdbfsf dfsndjdf sd sbdf fwbefiubw bfiwbfib lorem ipsum sdfjs fdsdj fshdbfsf dfsndjdf sd sbdf fwbefiubw bfiwbfib lorem ipsum sdfjs fdsdj fshdbfsf dfsndjdf sd sbdf fwbefiubw bfiwbfib
	
\section{Tipo o Nivel de Investigaci\'on}

	lorem ipsum sdfjs fdsdj fshdbfsf dfsndjdf sd sbdf fwbefiubw bfiwbfib lorem ipsum sdfjs fdsdj fshdbfsf dfsndjdf sd sbdf fwbefiubw bfiwbfib lorem ipsum sdfjs fdsdj fshdbfsf dfsndjdf sd sbdf fwbefiubw bfiwbfib lorem ipsum sdfjs fdsdj fshdbfsf dfsndjdf sd sbdf fwbefiubw bfiwbfib lorem ipsum sdfjs fdsdj fshdbfsf dfsndjdf sd sbdf fwbefiubw bfiwbfib lorem ipsum sdfjs fdsdj fshdbfsf dfsndjdf sd sbdf fwbefiubw bfiwbfib lorem ipsum sdfjs fdsdj fshdbfsf dfsndjdf sd sbdf fwbefiubw bfiwbfib lorem ipsum sdfjs fdsdj fshdbfsf dfsndjdf sd sbdf fwbefiubw bfiwbfib	
	
\section{Dise\~no de Investigaci\'on}
	
	lorem ipsum sdfjs fdsdj fshdbfsf dfsndjdf sd sbdf fwbefiubw bfiwbfib lorem ipsum sdfjs fdsdj fshdbfsf dfsndjdf sd sbdf fwbefiubw bfiwbfib lorem ipsum sdfjs fdsdj fshdbfsf dfsndjdf sd sbdf fwbefiubw bfiwbfib lorem ipsum sdfjs fdsdj fshdbfsf dfsndjdf sd sbdf fwbefiubw bfiwbfib lorem ipsum sdfjs fdsdj fshdbfsf dfsndjdf sd sbdf fwbefiubw bfiwbfib lorem ipsum sdfjs fdsdj fshdbfsf dfsndjdf sd sbdf fwbefiubw bfiwbfib lorem ipsum sdfjs fdsdj fshdbfsf dfsndjdf sd sbdf fwbefiubw bfiwbfib lorem ipsum sdfjs fdsdj fshdbfsf dfsndjdf sd sbdf fwbefiubw bfiwbfib
	

\section{Metodolog\'ia}


	lorem ipsum sdfjs fdsdj fshdbfsf dfsndjdf sd sbdf fwbefiubw bfiwbfib lorem ipsum sdfjs fdsdj fshdbfsf dfsndjdf sd sbdf fwbefiubw bfiwbfib lorem ipsum sdfjs fdsdj fshdbfsf dfsndjdf sd sbdf fwbefiubw bfiwbfib lorem ipsum sdfjs fdsdj fshdbfsf dfsndjdf sd sbdf fwbefiubw bfiwbfib lorem ipsum sdfjs fdsdj fshdbfsf dfsndjdf sd sbdf fwbefiubw bfiwbfib lorem ipsum sdfjs fdsdj fshdbfsf dfsndjdf sd sbdf fwbefiubw bfiwbfib lorem ipsum sdfjs fdsdj fshdbfsf dfsndjdf sd sbdf fwbefiubw bfiwbfib lorem ipsum sdfjs fdsdj fshdbfsf dfsndjdf sd sbdf fwbefiubw bfiwbfib


	
%\include{conclu}

\appendix
%% Cap'itulos incluidos despues del comando \appendix aparecen como ap'endices
%% de la tesis.
%\include{apendiceA}
%\include{apendiceB}
%\include{apendiceC}

%% Incluir la bibliograf'ia. Mirar el archivo "biblio.bib" para m'as detales
%% y un ejemplo.
\bibliographystyle{apalike}
\bibliography{biblio}

\end{document}
